\section{Introduction}
\label{sec:intro}

The rapid proliferation of Large Language Models (LLMs) has fundamentally altered the landscape of natural language processing and enterprise knowledge management. While early monolithic models relied entirely on their parametric memory, they inherently suffered from static knowledge cutoffs and high hallucination rates. These limitations are increasingly mitigated by Retrieval-Augmented Generation (RAG) \cite{lewis2020rag}. RAG architectures bifurcate the intelligence of the system: preserving the LLM as a stateless reasoning engine while delegating factual recall to a dynamic, non-parametric external memory, typically instantiated as a dense vector database.

This architectural shift enables specialized AI agents capable of reasoning over domain-specific corpora, ranging from legal case files to electronic health records. However, the centralization of this sensitive non-parametric memory introduces new privacy considerations. Standard RAG implementations process and store corporate knowledge bases as dense embeddings within a single monolithic vector database. Recent advancements in embedding inversion attacks \cite{carlini2021extract, morris2023vec2text} have demonstrated that dense vectors can leak substantial information. If an adversary gains access to a centralized repository, or successfully mounts an extraction attack via a local LLM interface, sensitive semantic clusters can be compromised.

To mitigate this risk, organizations often partition the underlying knowledge base into multiple logically isolated storage units, or \textit{shards} \cite{wang2024mrag}. In a distributed deployment, the RAG orchestration layer routes a query only to relevant shards. The foundational premise of semantic sharding as a defense-in-depth mechanism is isolation: an adversary who successfully compromises a single shard should only observe fragmented, incoherent data, severely limiting the probability of full topic reconstruction. We emphasize that sharding is a software-layer mitigation, not a substitute for encryption or access control.

However, the design of the partitioning algorithm dictates the success of this isolation. Standard industry approaches rely on unsupervised semantic clustering, such as K-Means or modularity-based community detection \cite{ng2001spectral}, which aim to maximize intra-shard semantic coherence to improve retrieval utility. Our research indicates that these algorithms inadvertently co-locate all documents pertaining to a specific sensitive topic onto the same shard, maximizing single-shard disclosure risk. Conversely, privacy-first frameworks explicitly scramble the semantic layout to enforce strict data isolation \cite{yang2025splitrag}, but inherently destroy the semantic proximity required by dense retrievers, leading to significant degradations in standard machine-learning retrieval metrics (e.g., Recall@k, MRR). 

This paper introduces a machine-learning methodology for exposure-aware representation-space partitioning in retrieval-augmented systems. We empirically demonstrate the inherent trade-off between RAG retrieval utility and data isolation, proposing a tunable graph-based optimization framework. We introduce the SHArd Disclosure Estimator (SHADE), a graph-based proxy metric that estimates the risk of semantic co-location. Building upon this metric, we propose CoShard, a graph partitioning algorithm that integrates the SHADE penalty into its modularity maximization objective. By formally penalizing the co-location of highly related documents, CoShard forces the distribution of conceptual topics across multiple shards, reducing single-shard topological exposure.

The core contributions of this work are summarized as follows:

\begin{itemize}
    \item \textbf{Exposure-Aware Partitioning Objective:} We formalize the representation-space partitioning problem as a trade-off between semantic retrieval utility and intra-shard semantic disclosure risk.
    \item \textbf{The SHADE Proxy Metric:} We introduce SHADE, an efficiently computable heuristic that models the graph-theoretic concentration of sensitive semantic edges within a partition.
    \item \textbf{The CoShard Algorithm:} We propose a scalable community detection algorithm that natively incorporates the SHADE metric as an optimization constraint, allowing practitioners to tune the precise boundary between retrieval utility and privacy.
    \item \textbf{Reproducible Empirical Methodology:} We conduct rigorous experiments over two public, domain-specific proxy benchmarks (MedQA and LegalBench). We evaluate CoShard against established baselines using standard ML retrieval metrics (Recall@k, MRR, nDCG@10) alongside an independent Entity Leakage Rate metric, demonstrating that topology-aware sharding can significantly reduce proxy indicators of reconstruction exposure while preserving retrieval performance.
\end{itemize}

The remainder of this paper is structured as follows. Section \ref{sec:related_work} surveys the literature surrounding privacy-preserving LLMs. Section \ref{sec:methodology} details the mathematical formulation of the SHADE metric and the CoShard optimization objective. Section \ref{sec:experiments} presents the empirical evaluation methodology, the baseline comparisons, and an ablation study. Section \ref{sec:privacy_analysis} provides an interpretive risk model analysis of topological exposure. Section \ref{sec:discussion} critically discusses the implications and limitations of this methodology, and Section \ref{sec:conclusion} concludes.
